\documentclass[12pt, a4paper]{book}

% ---------- Packages ----------
\usepackage{amsmath, amssymb, amsthm}
\usepackage{mathtools}
\usepackage{geometry}
\usepackage{hyperref}
\usepackage{enumitem}
\usepackage{booktabs}
\usepackage{array}
\usepackage{xcolor}
\usepackage{tcolorbox}
\usepackage{fancyhdr}
\usepackage{titlesec}
\usepackage{microtype}
\tcbuselibrary{theorems, skins, breakable}

% ---------- Page Geometry ----------
\geometry{top=2.5cm, bottom=2.5cm, left=3cm, right=3cm}

% ---------- Header/Footer ----------
\pagestyle{fancy}
\fancyhf{}
\fancyhead[LE,RO]{\thepage}
\fancyhead[RE]{\textit{A-Level Further Mathematics}}
\fancyhead[LO]{\textit{\leftmark}}
\renewcommand{\headrulewidth}{0.4pt}

% ---------- Theorem Environments ----------
\tcbset{
  thmstyle/.style={
    enhanced, breakable,
    colback=gray!5, colframe=black!60,
    fonttitle=\bfseries, separator sign={ },
    left=4pt, right=4pt, top=4pt, bottom=4pt,
  },
  defstyle/.style={
    enhanced, breakable,
    colback=blue!4, colframe=blue!50!black,
    fonttitle=\bfseries,
    left=4pt, right=4pt, top=4pt, bottom=4pt,
  },
  exstyle/.style={
    enhanced, breakable,
    colback=green!3, colframe=green!50!black,
    fonttitle=\bfseries,
    left=4pt, right=4pt, top=4pt, bottom=4pt,
  },
  keystyle/.style={
    enhanced,
    colback=yellow!10, colframe=orange!70!black,
    fonttitle=\bfseries,
    left=4pt, right=4pt, top=3pt, bottom=3pt,
  },
}

\newtcbtheorem[number within=section]{theorem}{Theorem}{thmstyle}{thm}
\newtcbtheorem[number within=section]{definition}{Definition}{defstyle}{def}
\newtcbtheorem[number within=section]{example}{Example}{exstyle}{ex}
\newtcbtheorem[number within=section]{proposition}{Proposition}{thmstyle}{prop}
\newtcbtheorem[number within=section]{corollary}{Corollary}{thmstyle}{cor}

\newenvironment{keyresult}{%
  \begin{tcolorbox}[keystyle]%
}{%
  \end{tcolorbox}%
}

% ---------- Math Operators ----------
\DeclareMathOperator{\sech}{sech}
\DeclareMathOperator{\csch}{csch}
\DeclareMathOperator{\arcsinh}{arcsinh}
\DeclareMathOperator{\arccosh}{arccosh}
\DeclareMathOperator{\arctanh}{arctanh}
\newcommand{\arsinh}{\operatorname{arsinh}}
\newcommand{\arcosh}{\operatorname{arcosh}}
\newcommand{\artanh}{\operatorname{artanh}}
\newcommand{\R}{\mathbb{R}}
\newcommand{\C}{\mathbb{C}}
\newcommand{\Z}{\mathbb{Z}}
\newcommand{\N}{\mathbb{N}}
\newcommand{\dd}{\,\mathrm{d}}
\newcommand{\dx}{\,\mathrm{d}x}
\newcommand{\dt}{\,\mathrm{d}t}
\newcommand{\ds}{\,\mathrm{d}s}
\newcommand{\e}{\mathrm{e}}
\newcommand{\ii}{\mathrm{i}}
\newcommand{\trans}{^{\mathsf{T}}}
\newcommand{\inv}{^{-1}}
\renewcommand{\vec}[1]{\mathbf{#1}}

% ---------- Title ----------
\title{%
  \Huge\textbf{A-Level Further Mathematics}\\[0.5em]
  \Large Cambridge Preparation Notes\\[1em]
  \large Pure Mathematics: Hyperbolic Functions, Matrices,\\
  Differentiation, Integration, Complex Numbers,\\
  Differential Equations
}
\author{Rayyan Hanif bin Rizal, Kolej MARA Seremban}
\date{}

% =========================================================
\begin{document}
% =========================================================

\maketitle
\tableofcontents
\newpage

% =========================================================
\chapter{Hyperbolic Functions}
% =========================================================

\section{Definitions and Exponential Forms}

\begin{definition}{Hyperbolic Functions}{}
For $x \in \R$:
\[
\sinh x = \frac{\e^x - \e^{-x}}{2}, \qquad
\cosh x = \frac{\e^x + \e^{-x}}{2}, \qquad
\tanh x = \frac{\sinh x}{\cosh x} = \frac{\e^x - \e^{-x}}{\e^x + \e^{-x}}.
\]
The reciprocal functions are $\csch x = 1/\sinh x$, $\sech x = 1/\cosh x$, $\coth x = 1/\tanh x$.
\end{definition}

\textit{Domain and range:} $\sinh:\R\to\R$ (bijection); $\cosh:\R\to[1,\infty)$ (even, minimum at $x=0$); $\tanh:\R\to(-1,1)$.

\section{Fundamental Identities}

\begin{keyresult}
\textbf{Osborn's Rule.} Hyperbolic identities are obtained from trigonometric identities by replacing $\sin\theta$ with $\sinh x$ and $\cos\theta$ with $\cosh x$, \emph{and changing the sign of every product of two sinh terms}.
\end{keyresult}

Key identities (derivable directly from definitions):
\begin{align}
\cosh^2 x - \sinh^2 x &= 1 \label{eq:hpythagorean}\\
\sinh(x \pm y) &= \sinh x\cosh y \pm \cosh x\sinh y \\
\cosh(x \pm y) &= \cosh x\cosh y \pm \sinh x\sinh y \\
\sinh 2x &= 2\sinh x\cosh x \\
\cosh 2x &= \cosh^2 x + \sinh^2 x = 2\cosh^2 x - 1 = 1 + 2\sinh^2 x \\
1 - \tanh^2 x &= \sech^2 x, \qquad \coth^2 x - 1 = \csch^2 x.
\end{align}

\textit{Derivation of \eqref{eq:hpythagorean}:}
\[
\cosh^2 x - \sinh^2 x
= \frac{(\e^x+\e^{-x})^2 - (\e^x-\e^{-x})^2}{4}
= \frac{4}{4} = 1.
\]

\section{Differentiation and Integration}

\begin{proposition}{Derivatives of Hyperbolic Functions}{}
\[
\frac{\dd}{\dx}\sinh x = \cosh x, \quad
\frac{\dd}{\dx}\cosh x = \sinh x, \quad
\frac{\dd}{\dx}\tanh x = \sech^2 x.
\]
\end{proposition}

\textit{Proof:} Differentiate the exponential definitions term by term. $\frac{\dd}{\dx}\sinh x = \frac{\e^x+\e^{-x}}{2}=\cosh x$, and similarly for others.

Standard integrals follow immediately:
\[
\int\sinh x\dx = \cosh x + C,\quad
\int\cosh x\dx = \sinh x + C,\quad
\int\sech^2 x\dx = \tanh x + C.
\]

\section{Inverse Hyperbolic Functions}

\begin{theorem}{Logarithmic Forms}{}
\begin{align*}
\arsinh x &= \ln\!\left(x + \sqrt{x^2+1}\right), \quad x\in\R,\\
\arcosh x &= \ln\!\left(x + \sqrt{x^2-1}\right), \quad x\geq 1,\\
\artanh x &= \tfrac{1}{2}\ln\!\left(\frac{1+x}{1-x}\right), \quad |x|<1.
\end{align*}
\end{theorem}

\textit{Derivation of} $\arsinh$: Let $y=\sinh\inv x$, so $x=\sinh y=\frac{\e^y-\e^{-y}}{2}$. Then $\e^{2y}-2x\e^y-1=0$. Solving (positive root):
\[
\e^y = x+\sqrt{x^2+1} \implies y = \ln(x+\sqrt{x^2+1}).
\]

\begin{proposition}{Derivatives of Inverse Hyperbolic Functions}{}
\[
\begin{aligned}
\frac{\mathrm{d}}{\mathrm{d}x}\arsinh x &= \frac{1}{\sqrt{x^2+1}}, \\
\frac{\mathrm{d}}{\mathrm{d}x}\arcosh x &= \frac{1}{\sqrt{x^2-1}} \quad (x>1), \\
\frac{\mathrm{d}}{\mathrm{d}x}\artanh x &= \frac{1}{1-x^2} \quad (|x|<1).
\end{aligned}
\]
\end{proposition}

\textit{Derivation:} Implicit differentiation. For $\displaystyle y=\arsinh x$: $\displaystyle x=\sinh y$, so $\displaystyle \frac{\dx}{\dd y}=\cosh y=\sqrt{1+\sinh^2 y}=\sqrt{1+x^2}$, giving $\displaystyle \frac{\dd y}{\dx}=\frac{1}{\sqrt{1+x^2}}$.

\begin{keyresult}
\textbf{Standard integral forms:}
\[
\begin{aligned}
\int\frac{\dx}{\sqrt{x^2+a^2}} = \arsinh\frac{x}{a}+C = \ln\!\left(x+\sqrt{x^2+a^2}\right)+C, \\
\int\frac{\dx}{\sqrt{x^2-a^2}} = \arcosh\frac{x}{a}+C\quad(x>a>0), \\
\int\frac{\dx}{a^2-x^2} = \frac{1}{a}\artanh\frac{x}{a}+C\quad(|x|<a).
\end{aligned}
\]
\end{keyresult}

\begin{example}{Inverse Hyperbolic Integral}{}
Evaluate $\displaystyle \int_0^2 \frac{\mathrm{d}x}{\sqrt{4+x^2}}$.

\medskip

\textbf{Solution.} Using $\displaystyle \text{arsinh}(x/2)$:
\[
\left[\arsinh\frac{x}{2}\right]_0^2 = \arsinh 1 - \arsinh 0 = \ln(1+\sqrt{2}).
\]
\end{example}

\section*{Exercises — Chapter 1}
\begin{enumerate}[label=\textbf{\arabic*.}]
  \item Show that $\displaystyle \cosh 2x = 1 + 2\sinh^2 x$ using the exponential definition.
  \item Solve $\displaystyle 2\cosh x + \sinh x = 3$. (Express your answer as a natural logarithm.)
  \item Differentiate $\displaystyle y = \arsinh(2x+1)$ and simplify fully.
  \item Evaluate $\displaystyle\int_0^1\tanh x\dx$. (\textit{Hint:} $\tanh x = \frac{\sinh x}{\cosh x}$.)
  \item \textbf{[Extension]} Prove that $\artanh(\tanh 2x) = 2x$ for all $x\in\R$, and use the logarithmic form of $\artanh$ to express $\tanh 2x$ as a ratio of exponentials. Verify using the definition.
\end{enumerate}

% =========================================================
\chapter{Matrices 2: Eigenvalues, Eigenvectors, and Diagonalisation}
% =========================================================

\section{Eigenvalues and Eigenvectors}

\begin{definition}{Eigenvalue and Eigenvector}{}
Let $A$ be an $n\times n$ matrix. A scalar $\lambda\in\C$ is an \emph{eigenvalue} of $A$ if there exists a non-zero vector $\vec{v}\in\C^n$ such that
\[
A\vec{v} = \lambda\vec{v}.
\]
Any such $\vec{v}$ is an \emph{eigenvector} corresponding to $\lambda$.
\end{definition}

\textbf{Finding eigenvalues.} The equation $A\vec{v}=\lambda\vec{v}$ has a non-trivial solution iff
\[
\det(A - \lambda I) = 0.
\]
This is the \emph{characteristic equation}. The polynomial $\chi_A(\lambda)=\det(A-\lambda I)$ is the \emph{characteristic polynomial} of degree $n$.

\textbf{Finding eigenvectors.} For each eigenvalue $\lambda_i$, solve $(A-\lambda_i I)\vec{v}=\vec{0}$ (row-reduce). The solution space is the \emph{eigenspace} $E_{\lambda_i}$.

\begin{example}{$3\times 3$ Eigenvalue Problem}{}
Find the eigenvalues and eigenvectors of
$A = \begin{pmatrix}4&1&0\\2&3&0\\0&0&5\end{pmatrix}$.

\textit{Characteristic equation:}
\[
\det(A-\lambda I)=\det\begin{pmatrix}4-\lambda&1&0\\2&3-\lambda&0\\0&0&5-\lambda\end{pmatrix}
=(5-\lambda)\bigl[(4-\lambda)(3-\lambda)-2\bigr]=0.
\]
Inner factor: $(4-\lambda)(3-\lambda)-2=\lambda^2-7\lambda+10=(\lambda-2)(\lambda-5)$.

Eigenvalues: $\lambda_1=2,\;\lambda_2=5$ (double).

For $\lambda=2$: $(A-2I)\vec{v}=\vec{0}$ gives $\begin{pmatrix}2&1&0\\2&1&0\\0&0&3\end{pmatrix}\vec{v}=\vec{0}$, so $v_3=0$, $2v_1+v_2=0$. Eigenvector: $\begin{pmatrix}1\\-2\\0\end{pmatrix}$.

For $\lambda=5$: $(A-5I)\vec{v}=\vec{0}$ gives $\begin{pmatrix}-1&1&0\\2&-2&0\\0&0&0\end{pmatrix}$, so $v_1=v_2$, $v_3$ free. Two independent eigenvectors: $\begin{pmatrix}1\\1\\0\end{pmatrix}$, $\begin{pmatrix}0\\0\\1\end{pmatrix}$.
\end{example}

\section{Diagonalisation}

\begin{theorem}{Diagonalisation}{}
An $n\times n$ matrix $A$ is \emph{diagonalisable} iff it has $n$ linearly independent eigenvectors. If $P$ is the matrix whose columns are these eigenvectors, then
\[
P\inv A P = D = \operatorname{diag}(\lambda_1,\ldots,\lambda_n).
\]
Equivalently, $A = PDP\inv$.
\end{theorem}

\textit{Proof sketch.} The $k$-th column of $AP$ equals $A\vec{p}_k=\lambda_k\vec{p}_k$, which is the $k$-th column of $PD$. Thus $AP=PD$, and invertibility of $P$ (guaranteed by independence) gives $P\inv AP=D$.

\begin{keyresult}
\textbf{Matrix powers:} $A^n = PD^n P\inv$, where $D^n=\operatorname{diag}(\lambda_1^n,\ldots,\lambda_k^n)$.
\end{keyresult}

\begin{example}{Computing $A^n$ via Diagonalisation}{}
Let $A=\begin{pmatrix}3&1\\1&3\end{pmatrix}$. Eigenvalues: $\lambda=2,4$ with eigenvectors $\begin{pmatrix}1\\-1\end{pmatrix},\begin{pmatrix}1\\1\end{pmatrix}$.

\medskip

With the modal matrix $P$ and its inverse $P^{-1}$ found as:
\[
P=\begin{pmatrix}1&1\\-1&1\end{pmatrix}, \quad P^{-1}=\tfrac{1}{2}\begin{pmatrix}1&-1\\1&1\end{pmatrix}
\]
We can compute $A^n$ by splitting the calculation across lines:
\begin{align*}
A^n &= P\begin{pmatrix}2^n&0\\0&4^n\end{pmatrix}P^{-1} \\
&= \tfrac{1}{2}\begin{pmatrix}2^n+4^n & 4^n-2^n \\ 4^n-2^n & 2^n+4^n\end{pmatrix}.
\end{align*}
\end{example}

\section{Symmetric Matrices and Orthogonal Diagonalisation}

\begin{theorem}{Spectral Theorem (Real Symmetric Matrices)}{}
If $A$ is a real symmetric matrix ($A=A\trans$), then:
\begin{enumerate}[label=(\roman*)]
  \item All eigenvalues of $A$ are real.
  \item Eigenvectors corresponding to distinct eigenvalues are orthogonal.
  \item $A$ is orthogonally diagonalisable: $A=QDQ\trans$ where $Q$ is orthogonal ($QQ\trans=I$).
\end{enumerate}
\end{theorem}

\textit{Proof (reality of eigenvalues):} Let $A\vec{v}=\lambda\vec{v}$, $\vec{v}\neq\vec{0}$. Then $\bar{\vec{v}}\trans A\vec{v}=\lambda\|\vec{v}\|^2$. Since $A=A\trans=\bar{A}$, the left side is also $\overline{\bar{\vec{v}}\trans A\vec{v}}=\bar{\lambda}\|\vec{v}\|^2$. Hence $\lambda=\bar{\lambda}\in\R$.

\section{Systems of Differential Equations via Matrices}

A linear system $\dot{\vec{x}}=A\vec{x}$ (where $A$ is diagonalisable with $A=PDP\inv$) is decoupled by the substitution $\vec{x}=P\vec{y}$, giving $\dot{\vec{y}}=D\vec{y}$, i.e.\ $\dot{y}_k=\lambda_k y_k$ for each $k$. Thus $y_k=C_k\e^{\lambda_k t}$ and $\vec{x}=P\vec{y}$.

\begin{example}{Decoupled System}{}
Solve $\dot{x}=3x+y,\;\dot{y}=x+3y$ with $x(0)=1,y(0)=0$.

From the previous example, $\lambda=2,4$, $P=\begin{pmatrix}1&1\\-1&1\end{pmatrix}$. The general solution is:
\[
\begin{pmatrix}x\\y\end{pmatrix}=C_1\e^{2t}\begin{pmatrix}1\\-1\end{pmatrix}+C_2\e^{4t}\begin{pmatrix}1\\1\end{pmatrix}.
\]
Initial conditions: $C_1+C_2=1$, $-C_1+C_2=0\implies C_1=C_2=\tfrac{1}{2}$.
\[
x=\tfrac{1}{2}(\e^{2t}+\e^{4t}),\quad y=\tfrac{1}{2}(\e^{4t}-\e^{2t}).
\]
\end{example}

\section*{Exercises --- Chapter 2}
\begin{enumerate}[label=\textbf{\arabic*.}]
  \item Find the eigenvalues and eigenvectors of $\begin{pmatrix}1&2\\3&2\end{pmatrix}$.
  \item Show that $A=\begin{pmatrix}0&1\\-1&2\end{pmatrix}$ is not diagonalisable over $\R$.
  \item If $A$ has eigenvalues $\lambda_1,\lambda_2$, show $\det A=\lambda_1\lambda_2$ and $\operatorname{tr}A=\lambda_1+\lambda_2$.
  \item Find $A^{10}$ for $A=\begin{pmatrix}5&3\\-2&0\end{pmatrix}$ using diagonalisation.
  \item \textbf{[Extension]} Let $A=\begin{pmatrix}2&1&0\\0&2&1\\0&0&3\end{pmatrix}$. Determine whether $A$ is diagonalisable. If not, explain what obstruction arises. Find $A^n$ for the $3\times 3$ block by writing $A=D+N$ where $D$ is diagonal and $N$ is nilpotent, and use the fact that $DN=ND$ in this case.
\end{enumerate}

% =========================================================
\chapter{Differentiation}
% =========================================================

\section{Implicit Differentiation}

If $\displaystyle F(x,y)=0$ defines $\displaystyle y$ implicitly as a function of $\displaystyle x$, differentiate with respect to $\displaystyle x$ using the chain rule: wherever $\displaystyle y$ appears, multiply by $\displaystyle \frac{\dd y}{\dx}$.

\begin{example}{Implicit Differentiation}{}
Find $\displaystyle \frac{\dd y}{\dx}$ for $\displaystyle x^3 + y^3 = 3xy$ (the Folium of Descartes).

Differentiating: $\displaystyle 3x^2 + 3y^2\frac{\dd y}{\dx} = 3y + 3x\frac{\dd y}{\dx}$.

\[
\frac{\dd y}{\dx}(3y^2-3x) = 3y-3x^2 \implies \frac{\dd y}{\dx}=\frac{y-x^2}{y^2-x}.
\]
\end{example}

\textbf{Second derivatives implicitly.} Differentiate $\displaystyle \frac{\dd y}{\dx}=f(x,y)$ with respect to $\displaystyle x$, again using the chain rule on $y$ terms.

\section{Parametric Differentiation}

For a curve $\displaystyle x=f(t),\;y=g(t)$:
\begin{align}
\frac{\dd y}{\dx} &= \frac{\dot{g}(t)}{\dot{f}(t)}, \label{eq:param1}\\
\frac{\dd^2 y}{\dx^2} &= \frac{\dfrac{\dd}{\dt}\!\left(\dfrac{\dd y}{\dx}\right)}{\dot{f}(t)}. \label{eq:param2}
\end{align}

\begin{example}{Parametric Second Derivative}{}
For $\displaystyle x=t^2,\;y=t^3$: $\frac{\dd y}{\dx}=\frac{3t^2}{2t}=\frac{3t}{2}$.
\[
\frac{\dd^2 y}{\dx^2}=\frac{\frac{\dd}{\dt}(\frac{3t}{2})}{2t}=\frac{3/2}{2t}=\frac{3}{4t}.
\]
\end{example}

\begin{keyresult}
\textbf{Arc length (parametric):}
\[
L = \int_{t_1}^{t_2}\sqrt{\left(\frac{\dx}{\dt}\right)^2+\left(\frac{\dd y}{\dt}\right)^2}\dt.
\]
\end{keyresult}

\section{Maclaurin and Taylor Series}

\begin{definition}{Taylor and Maclaurin Series}{}
If $\displaystyle f$ is infinitely differentiable at $a$:
\[
f(x) = \sum_{n=0}^{\infty}\frac{f^{(n)}(a)}{n!}(x-a)^n.
\]
For $\displaystyle a=0$ this is the \emph{Maclaurin series} of $f$.
\end{definition}

\begin{keyresult}
\textbf{Standard Maclaurin series} (valid for indicated domains):
\begin{align*}
\e^x &= \sum_{n=0}^{\infty}\frac{x^n}{n!} = 1+x+\frac{x^2}{2!}+\cdots & (x\in\R),\\
\sin x &= \sum_{n=0}^{\infty}\frac{(-1)^n x^{2n+1}}{(2n+1)!} & (x\in\R),\\
\cos x &= \sum_{n=0}^{\infty}\frac{(-1)^n x^{2n}}{(2n)!} & (x\in\R),\\
\ln(1+x) &= \sum_{n=1}^{\infty}\frac{(-1)^{n+1}x^n}{n} = x-\frac{x^2}{2}+\frac{x^3}{3}-\cdots & (-1<x\leq 1),\\
(1+x)^k &= \sum_{n=0}^{\infty}\binom{k}{n}x^n & (|x|<1),\\
\sinh x &= x+\frac{x^3}{3!}+\frac{x^5}{5!}+\cdots & (x\in\R),\\
\cosh x &= 1+\frac{x^2}{2!}+\frac{x^4}{4!}+\cdots & (x\in\R).
\end{align*}
\end{keyresult}

\textbf{Combining series.} New series are obtained by substitution, differentiation, or integration of known series.

\begin{example}{Series by Composition}{}
Find the Maclaurin series for $\ln(\cos x)$ up to the $x^4$ term.

$\displaystyle \cos x = 1 - \frac{x^2}{2}+\frac{x^4}{24}-\cdots$. Let $\displaystyle u=-\frac{x^2}{2}+\frac{x^4}{24}-\cdots$ so $\displaystyle \ln(\cos x)=\ln(1+u)=u-\frac{u^2}{2}+\cdots$
\[
= -\frac{x^2}{2}+\frac{x^4}{24} - \frac{1}{2}\!\left(\frac{x^2}{2}\right)^2+O(x^6)
= -\frac{x^2}{2}-\frac{x^4}{12}+O(x^6).
\]
\end{example}

\section{L'H\^opital's Rule}

\begin{theorem}{L'H\^opital's Rule}{}
If $\displaystyle \lim_{x\to a}f(x)=\lim_{x\to a}g(x)=0$ (or both $\displaystyle \pm\infty$) and $\displaystyle g'(x)\neq0$ near $\displaystyle a$, then
\[
\lim_{x\to a}\frac{f(x)}{g(x)}=\lim_{x\to a}\frac{f'(x)}{g'(x)},
\]
provided the right-hand limit exists.
\end{theorem}

The rule applies repeatedly. Indeterminate forms $\displaystyle 0\cdot\infty$, $\displaystyle \infty-\infty$, $\displaystyle 0^0$, $\displaystyle \infty^0$, $\displaystyle 1^\infty$ are first converted to $\displaystyle 0/0$ or $\displaystyle \infty/\infty$.

\section*{Exercises --- Chapter 3}
\begin{enumerate}[label=\textbf{\arabic*.}]
  \item Find $\displaystyle \dfrac{\dd y}{\dx}$ given $\displaystyle \e^{xy}+y^2=x$.
  \item For $\displaystyle x=a\cos t,\;y=b\sin t$, find $\displaystyle \dfrac{\dd^2 y}{\dx^2}$ in terms of $\displaystyle t$.
  \item Find the Maclaurin series for $\displaystyle \tan x$ up to and including the $\displaystyle x^5$ term. (\textit{Hint:} use $\displaystyle \sin x = \cos x\displaystyle \cdot\tan x$ and compare series.)
  \item Use L'H\^opital's rule to evaluate $\displaystyle \lim_{x\to 0}\dfrac{x-\sin x}{x^3}$.
  \item \textbf{[Extension]} Derive the Maclaurin series for $\displaystyle \arctan x$ by integrating $\displaystyle \frac{1}{1+x^2}$ term by term. Hence show $\displaystyle \frac{\pi}{4}=1-\frac{1}{3}+\frac{1}{5}-\cdots$ (Leibniz formula), and estimate the error in approximating $\displaystyle \pi/4$ by the partial sum of $\displaystyle n$ terms.
\end{enumerate}

% =========================================================
\chapter{Integration}
% =========================================================

\section{Integration by Parts}

\begin{keyresult}
$\displaystyle\int u\,\dd v = uv - \int v\,\dd u$, or equivalently $\displaystyle\int u\,v'\dx = uv - \int u'v\dx$.

\textbf{LIATE priority for $u$:} Logarithms, Inverse trig, Algebraic, Trigonometric, Exponential.
\end{keyresult}

\begin{example}{Repeated Integration by Parts}{}
$\displaystyle I=\int x^2\e^x\dx$. Apply twice:
\[
I = x^2\e^x - 2\int x\e^x\dx = x^2\e^x - 2(x\e^x-\e^x)+C = \e^x(x^2-2x+2)+C.
\]
\end{example}

\section{Integration by Substitution}

Standard substitutions:
\begin{center}
\begin{tabular}{lll}
\toprule
Integrand contains & Substitution & Result\\
\midrule
$\sqrt{a^2-x^2}$ & $x=a\sin\theta$ & Removes square root\\
$\sqrt{a^2+x^2}$ & $x=a\tan\theta$ or $x=a\sinh t$ & \\
$\sqrt{x^2-a^2}$ & $x=a\sec\theta$ or $x=a\cosh t$ & \\
\bottomrule
\end{tabular}
\end{center}

\begin{example}{Trigonometric Substitution}{}
$\displaystyle\int\frac{\dx}{(a^2-x^2)^{3/2}}$, let $x=a\sin\theta$:
\[
= \int\frac{a\cos\theta\dd\theta}{a^3\cos^3\theta}=\frac{1}{a^2}\int\sec^2\theta\dd\theta=\frac{\tan\theta}{a^2}+C=\frac{x}{a^2\sqrt{a^2-x^2}}+C.
\]
\end{example}

\section{Partial Fractions}

For a rational function $P(x)/Q(x)$ with $\deg P<\deg Q$, decompose $Q$ into factors:

\begin{center}
\begin{tabular}{ll}
\toprule
Factor in $Q(x)$ & Partial fraction form\\
\midrule
$(x-a)$ & $\dfrac{A}{x-a}$\\[6pt]
$(x-a)^n$ & $\dfrac{A_1}{x-a}+\cdots+\dfrac{A_n}{(x-a)^n}$\\[6pt]
$(x^2+bx+c)$ irreducible & $\dfrac{Ax+B}{x^2+bx+c}$\\
\bottomrule
\end{tabular}
\end{center}

\section{Reduction Formulae}

A \emph{reduction formula} expresses $I_n$ in terms of $I_{n-1}$ (or $I_{n-2}$), enabling computation for general $n$.

\begin{example}{Derivation: $I_n=\int_0^{\pi/2}\sin^n x\dx$}{}
Integrate by parts with $u=\sin^{n-1}x$, $\dd v=\sin x\dx$:
\[
I_n = \bigl[-\sin^{n-1}x\cos x\bigr]_0^{\pi/2}+(n-1)\int_0^{\pi/2}\sin^{n-2}x\cos^2 x\dx.
\]
The boundary term vanishes. Using $\cos^2 x=1-\sin^2 x$:
\[
I_n=(n-1)I_{n-2}-(n-1)I_n \implies I_n = \frac{n-1}{n}I_{n-2}.
\]
With $\displaystyle I_0=\pi/2$ and $\displaystyle I_1=1$, this gives Wallis-type results.
\end{example}

\begin{keyresult}
\textbf{Wallis's formula:}
\[
\frac{\pi}{2}=\prod_{n=1}^{\infty}\frac{(2n)(2n)}{(2n-1)(2n+1)}=\frac{2\cdot2}{1\cdot3}\cdot\frac{4\cdot4}{3\cdot5}\cdot\frac{6\cdot6}{5\cdot7}\cdots
\]
\end{keyresult}

\section{Applications of Integration}

\subsection{Arc Length}

For $\displaystyle y=f(x)$ on $\displaystyle [a,b]$:
\[
L = \int_a^b\sqrt{1+\left(\frac{\dd y}{\dx}\right)^2}\dx.
\]

\subsection{Surface of Revolution}

Rotating $\displaystyle y=f(x)\geq 0$ about the $\displaystyle x$-axis on $\displaystyle [a,b]$:
\[
S = 2\pi\int_a^b y\sqrt{1+\left(\frac{\dd y}{\dx}\right)^2}\dx.
\]

\subsection{Volumes of Revolution}

Rotating $\displaystyle y=f(x)$ about the $\displaystyle x$-axis (disc method):
\[
V = \pi\int_a^b y^2\dx.
\]

\begin{example}{Volume via Disc Method}{}
Volume of sphere radius $\displaystyle r$: rotate $\displaystyle y=\sqrt{r^2-x^2}$ about the $\displaystyle x$-axis on $\displaystyle [-r,r]$:
\[
V=\pi\int_{-r}^{r}(r^2-x^2)\dx=\pi\left[r^2x-\frac{x^3}{3}\right]_{-r}^{r}=\frac{4\pi r^3}{3}.
\]
\end{example}

\section{Improper Integrals}

An integral is \emph{improper} if the interval is unbounded or the integrand is unbounded on the interval. Evaluate as a limit:
\[
\int_a^{\infty}f(x)\dx = \lim_{R\to\infty}\int_a^R f(x)\dx,
\]
converges if the limit exists and is finite.

\begin{example}{Gamma Function at $n=1$}{}
$\displaystyle\int_0^{\infty}x\e^{-x}\dx$. Integration by parts:
$\displaystyle \bigl[-x\e^{-x}\bigr]_0^{\infty}+\int_0^{\infty}\e^{-x}\dx = 0+1=1.$
\end{example}

\section*{Exercises --- Chapter 4}
\begin{enumerate}[label=\textbf{\arabic*.}]
  \item Find $\displaystyle\int x^3\ln x\dx$.
  \item Evaluate $\displaystyle\int\frac{3x+1}{(x+1)(x^2+1)}\dx$ using partial fractions.
  \item Establish the reduction formula $J_n = \displaystyle\int_0^1 x^n\e^x\dx = \e-nJ_{n-1}$, and compute $J_2$.
  \item Find the arc length of $y=\cosh x$ for $0\leq x\leq 1$.
  \item \textbf{[Extension]} Show that $\displaystyle\int_0^\infty\e^{-x^2}\dx=\frac{\sqrt{\pi}}{2}$ using the following argument: let $\displaystyle I=\int_0^\infty\e^{-x^2}\dx$, write $\displaystyle I^2=\int_0^\infty\int_0^\infty\e^{-(x^2+y^2)}\dx\dd y$, convert to polar coordinates, and evaluate.
\end{enumerate}

% =========================================================
\chapter{Complex Numbers}
% =========================================================

\section{Modulus-Argument Form and Euler's Formula}

\begin{keyresult}
\textbf{Euler's formula:} $\e^{\ii\theta}=\cos\theta+\ii\sin\theta$.

Every non-zero $\displaystyle z\in\C$ can be written $\displaystyle z=r\e^{\ii\theta}=r(\cos\theta+\ii\sin\theta)$, where $\displaystyle r=|z|>0$ and $\displaystyle \theta=\arg z$.
\end{keyresult}

\textit{Derivation from Maclaurin series:} Compare series expansions:
\[
\e^{\ii\theta}=1+\ii\theta+\frac{(\ii\theta)^2}{2!}+\frac{(\ii\theta)^3}{3!}+\cdots=\left(1-\frac{\theta^2}{2!}+\cdots\right)+\ii\left(\theta-\frac{\theta^3}{3!}+\cdots\right)=\cos\theta+\ii\sin\theta.
\]

\textbf{Multiplication and division:}
\[
z_1z_2 = r_1r_2\,\e^{\ii(\theta_1+\theta_2)}, \qquad
\frac{z_1}{z_2}=\frac{r_1}{r_2}\,\e^{\ii(\theta_1-\theta_2)}.
\]

\section{De Moivre's Theorem}

\begin{theorem}{De Moivre's Theorem}{}
For $\displaystyle n\in\Z$ (and more generally $\displaystyle n\in\R$ for principal values):
\[
(\cos\theta+\ii\sin\theta)^n = \cos n\theta+\ii\sin n\theta.
\]
\end{theorem}

\textit{Proof (integer $\displaystyle n$):} By Euler's formula, $\displaystyle (\e^{\ii\theta})^n=\e^{\ii n\theta}$.

\subsection{Multiple Angle Formulae}

From De Moivre with $\displaystyle z=\cos\theta+\ii\sin\theta$, let $\displaystyle z^n+z^{-n}=2\cos n\theta$ and $\displaystyle z^n-z^{-n}=2\ii\sin n\theta$.

Express powers of $\cos\theta$ and $\sin\theta$ in terms of multiples, or vice versa.

\begin{example}{Expressing $\cos^4\theta$}{}
Let $\displaystyle c=\cos\theta+\ii\sin\theta$. Then $\displaystyle 2\cos\theta=c+c^{-1}$, so
\[
(2\cos\theta)^4=(c+c^{-1})^4=c^4+4c^2+6+4c^{-2}+c^{-4}=2\cos4\theta+8\cos2\theta+6.
\]
Thus $\displaystyle \cos^4\theta=\frac{1}{8}(\cos4\theta+4\cos2\theta+3)$.
\end{example}

\section{Roots of Complex Numbers}

The equation $\displaystyle z^n=w$ has exactly $\displaystyle n$ roots. Writing $\displaystyle w=R\e^{\ii\phi}$:
\[
z_k = R^{1/n}\exp\!\left(\ii\,\frac{\phi+2k\pi}{n}\right), \quad k=0,1,\ldots,n-1.
\]
These $\displaystyle n$ roots lie equally spaced on a circle of radius $R^{1/n}$ in the Argand diagram.

\begin{example}{Cube Roots of Unity}{}
Solve $\displaystyle z^3=1$: $z_k=\e^{2k\pi\ii/3}$, $\displaystyle k=0,1,2$.
\[
z_0=1,\quad z_1=\e^{2\pi\ii/3}=-\tfrac{1}{2}+\tfrac{\sqrt{3}}{2}\ii\eqqcolon\omega,\quad z_2=\e^{4\pi\ii/3}=-\tfrac{1}{2}-\tfrac{\sqrt{3}}{2}\ii=\omega^2.
\]
\textbf{Key property:} $\displaystyle 1+\omega+\omega^2=0$, $\displaystyle \omega^3=1$.
\end{example}

\begin{keyresult}
\textbf{Sum of roots of unity:} For $\displaystyle \omega=\e^{2\pi\ii/n}$,
\[
\sum_{k=0}^{n-1}\omega^k = \frac{\omega^n-1}{\omega-1} = 0.
\]
\end{keyresult}

\section{Loci in the Argand Diagram}

Common loci arising from conditions on $z$:
\begin{itemize}
  \item $\displaystyle |z-a|=r$: circle centre $a$, radius $r$.
  \item $\displaystyle |z-a|=|z-b|$: perpendicular bisector of segment $ab$.
  \item $\displaystyle \arg(z-a)=\alpha$: ray from $a$ at angle $\alpha$ (excluding $a$).
  \item $\displaystyle |z-a|/|z-b|=k$ ($k\neq1$): circle (Apollonius circle).
  \item $\displaystyle \arg\!\left(\dfrac{z-a}{z-b}\right)=\alpha$: arc of circle through $\displaystyle a$ and $\displaystyle b$ subtending angle $\displaystyle \alpha$.
\end{itemize}

\begin{example}{Locus Problem}{}
Describe the locus of $z$ satisfying $\left|\dfrac{z-2}{z+2}\right|=2$.

This is $|z-2|=2|z+2|$, an Apollonius circle. Writing $z=x+\ii y$:
\[
(x-2)^2+y^2 = 4[(x+2)^2+y^2].
\]
Expanding and simplifying: $3x^2+3y^2+20x+12=0$, i.e.\ $(x+\tfrac{10}{3})^2+y^2=\tfrac{28}{9}$. Circle, centre $(-\tfrac{10}{3},0)$, radius $\tfrac{2\sqrt{7}}{3}$.
\end{example}

\section{Complex Exponentials and Trigonometry}

\begin{keyresult}
$\cos\theta=\dfrac{\e^{\ii\theta}+\e^{-\ii\theta}}{2}$, \quad
$\sin\theta=\dfrac{\e^{\ii\theta}-\e^{-\ii\theta}}{2\ii}$.

These unify hyperbolic and trigonometric functions: $\cos(\ii x)=\cosh x$, $\sin(\ii x)=\ii\sinh x$.
\end{keyresult}

\section*{Exercises --- Chapter 5}
\begin{enumerate}[label=\textbf{\arabic*.}]
  \item Use De Moivre's theorem to express $\sin 5\theta$ in terms of $\sin\theta$.
  \item Find all fifth roots of $-32$ in the form $r\e^{\ii\theta}$.
  \item Sketch the locus of $z$ such that $\arg(z-1-\ii)=\pi/4$, stating clearly which point is excluded.
  \item Show that $\displaystyle\sum_{k=0}^{n-1}\cos(a+kd)=\frac{\sin(nd/2)}{\sin(d/2)}\cos\!\left(a+\frac{(n-1)d}{2}\right)$ using the real part of a geometric series in $\e^{\ii d}$.
  \item \textbf{[Extension]} Let $\omega=\e^{2\pi\ii/7}$. Evaluate $\displaystyle\prod_{k=1}^{6}(2-\omega^k)$ by considering the polynomial $z^7-1$ factored over $\C$, and substituting $z=2$. Hence find $\displaystyle\prod_{k=1}^{6}|2-\omega^k|$.
\end{enumerate}

% =========================================================
\chapter{Differential Equations}
% =========================================================

\section{First-Order ODEs}

\subsection{Separable Equations}

An ODE of the form $\displaystyle \dfrac{\dd y}{\dx}=f(x)g(y)$ is separable:
\[
\int\frac{\dd y}{g(y)}=\int f(x)\dx+C.
\]

\subsection{Integrating Factor Method}

The linear first-order ODE $\displaystyle \dfrac{\dd y}{\dx}+P(x)y=Q(x)$ is solved by the integrating factor $\displaystyle \mu(x)=\e^{\int P(x)\dx}$:
\[
\frac{\dd}{\dx}(\mu y)=\mu Q \implies \mu y = \int\mu Q\dx + C.
\]

\textit{Derivation:} Multiply through by $\displaystyle \mu$: $\displaystyle \mu y'+\mu Py=\mu Q$. The left side is $\displaystyle (\mu y)'$ if $\displaystyle \mu'=\mu P$, i.e.\ $\displaystyle \mu=\e^{\int P\dx}$.

\begin{example}{Integrating Factor}{}
Solve $\displaystyle \dfrac{\dd y}{\dx}+\dfrac{2y}{x}=x^2$.

$\displaystyle P=2/x$, so $\displaystyle \mu=\e^{2\ln x}=x^2$. Multiply: $\displaystyle (x^2 y)'=x^4$.
\[
x^2 y = \frac{x^5}{5}+C \implies y=\frac{x^3}{5}+Cx^{-2}.
\]
\end{example}

\subsection{Homogeneous Equations}

$\displaystyle \dfrac{\dd y}{\dx}=f(y/x)$: substitute $v=y/x$ so $y=vx$, $y'=v+xv'$.

\subsection{Bernoulli Equations}

$\displaystyle \dfrac{\dd y}{\dx}+P(x)y=Q(x)y^n$: substitute $w=y^{1-n}$ to linearise.

\section{Second-Order Linear ODEs with Constant Coefficients}

The equation $\displaystyle ay''+by'+cy=f(x)$ has general solution $\displaystyle y=y_{\rm c}+y_{\rm p}$, where $\displaystyle y_{\rm c}$ (complementary function) solves the homogeneous equation and $\displaystyle _{\rm p}$ is any particular integral.

\subsection{Complementary Function}

Solve the \emph{auxiliary equation} $\displaystyle a\lambda^2+b\lambda+c=0$. Let $\displaystyle \Delta=b^2-4ac$.

\begin{keyresult}
\begin{center}
\begin{tabular}{lll}
\toprule
Case & Roots & Complementary function $\displaystyle y_{\rm c}$\\
\midrule
$\displaystyle \Delta>0$ & real distinct $\displaystyle \lambda_1\neq\lambda_2$ & $\displaystyle A\e^{\lambda_1 x}+B\e^{\lambda_2 x}$\\
$\displaystyle \Delta=0$ & repeated $\displaystyle \lambda$ & $\displaystyle (A+Bx)\e^{\lambda x}$\\
$\displaystyle \Delta<0$ & complex $\displaystyle \alpha\pm\ii\beta$ & $\displaystyle \e^{\alpha x}(A\cos\beta x+B\sin\beta x)$\\
\bottomrule
\end{tabular}
\end{center}
\end{keyresult}

\subsection{Particular Integral by Method of Undetermined Coefficients}

\begin{center}
\begin{tabular}{ll}
\toprule
$f(x)$ & Trial $y_{\rm p}$\\
\midrule
$P_n(x)$ (polynomial degree $n$) & $Q_n(x)$ (polynomial degree $n$)$^*$\\
$\e^{kx}$ & $C\e^{kx}$\;$^*$\\
$\e^{kx}P_n(x)$ & $\e^{kx}Q_n(x)$\;$^*$\\
$\cos kx$ or $\sin kx$ & $A\cos kx+B\sin kx$\;$^*$\\
\bottomrule
\end{tabular}
\end{center}
$^*$ Multiply by $x^s$ where $s$ is the smallest non-negative integer such that no term of the trial solution is in $y_{\rm c}$.

\begin{example}{Full Solution of Second-Order ODE}{}
Solve $y''-3y'+2y=\e^x$.

Auxiliary equation: $\lambda^2-3\lambda+2=(\lambda-1)(\lambda-2)=0$. Roots $\lambda=1,2$, so $y_{\rm c}=A\e^x+B\e^{2x}$.

Since $\e^x$ is in $y_{\rm c}$, try $y_{\rm p}=Cx\e^x$:
\[
y_{\rm p}'=C\e^x(1+x),\quad y_{\rm p}''=C\e^x(2+x).
\]
Substituting: $C\e^x(2+x)-3C\e^x(1+x)+2Cx\e^x=C\e^x(2+x-3-3x+2x)=C\e^x(x-1-x)\cdot... $

Let us redo: $y_{\rm p}''-3y_{\rm p}'+2y_{\rm p}=C\e^x(2+x)-3C\e^x(1+x)+2Cx\e^x=C\e^x[(2+x-3-3x+2x)]=C\e^x(2-3)=-C\e^x$.

So $-C\e^x=\e^x\implies C=-1$. General solution: $y=A\e^x+B\e^{2x}-x\e^x$.
\end{example}

\section{Variation of Parameters}

For $ay''+by'+cy=f(x)$ with complementary function $y_{\rm c}=Ay_1+By_2$, the particular integral is:
\[
y_{\rm p} = -y_1\int\frac{y_2 f}{W}\dx + y_2\int\frac{y_1 f}{W}\dx,
\]
where $W=y_1y_2'-y_2y_1'$ is the \emph{Wronskian}.

\section{Euler--Cauchy Equation}

The equation $ax^2y''+bxy'+cy=f(x)$ is solved by the substitution $x=\e^t$, transforming it to a constant-coefficient ODE in $t$.

\textit{Key identities under $x=\e^t$:}
\[
xy' = \dot{y},\qquad x^2y'' = \ddot{y}-\dot{y},
\]
where dots denote $\dd/\dt$.

\section{Simple Harmonic Motion and Damping}

The equation $\ddot{x}+\omega^2 x=0$ (simple harmonic motion) has solution $x=A\cos\omega t+B\sin\omega t$.

The damped oscillator $\ddot{x}+2k\dot{x}+\omega^2 x=0$:
\begin{itemize}
  \item $k<\omega$ (underdamped): $x=\e^{-kt}(A\cos\mu t+B\sin\mu t)$, $\mu=\sqrt{\omega^2-k^2}$.
  \item $k=\omega$ (critically damped): $x=(A+Bt)\e^{-kt}$.
  \item $k>\omega$ (overdamped): $x=A\e^{\lambda_1 t}+B\e^{\lambda_2 t}$, both $\lambda_i<0$.
\end{itemize}

\begin{example}{Forced Oscillator}{}
Solve $\ddot{x}+4x=3\cos t$, $x(0)=0$, $\dot{x}(0)=0$.

$y_{\rm c}=A\cos2t+B\sin2t$. Try $y_{\rm p}=C\cos t$: $-C\cos t+4C\cos t=3\cos t\implies C=1$.

So $x=A\cos2t+B\sin2t+\cos t$. Applying ICs: $A+1=0$, $2B=0$, giving $A=-1$, $B=0$.
\[
x=\cos t-\cos2t.
\]
\end{example}

\section{Resonance}

If $f(x)=\cos\omega t$ and the natural frequency is also $\omega$ (i.e.\ $\omega^2=\omega_{\rm n}^2$ in the auxiliary equation), the trial PI must be multiplied by $t$, giving \emph{resonance}: amplitude grows without bound as $t\to\infty$.

\section*{Exercises --- Chapter 6}
\begin{enumerate}[label=\textbf{\arabic*.}]
  \item Solve $\dfrac{\dd y}{\dx}=\dfrac{x(1+y^2)}{y(1+x^2)}$ with $y(0)=1$.
  \item Find the general solution of $y''-5y'+6y=\e^{2x}+x$.
  \item Solve the IVP $y''+2y'+5y=0$, $y(0)=1$, $y'(0)=0$.
  \item Use the substitution $v=y/x$ to solve $\dfrac{\dd y}{\dx}=\dfrac{y}{x}+\tan\!\left(\dfrac{y}{x}\right)$.
  \item \textbf{[Extension]} Consider $\ddot{x}+\omega^2 x=F\cos\omega t$ (resonant forcing). Assume the general solution takes the form $x=(A+Ct)\cos\omega t+(B+Dt)\sin\omega t$. Substitute into the ODE and show $D=F/(2\omega)$, $C=0$. With $x(0)=\dot{x}(0)=0$, write the full solution and describe the long-term behaviour.
\end{enumerate}

\newpage

\section*{Exploration Questions}

\hrulefill

\subsection*{Complex Numbers and De Moivre's Theorem}
\begin{itemize}
\item Investigate the geometric shape formed by the roots of $z^n = 1$ in the Argand diagram for different values of $n$. Describe how symmetry changes as $n$ increases.
\item Explore how the expression $(\cos \theta + i \sin \theta)^n$ can be used to generate trigonometric identities for $\cos(n\theta)$ and $\sin(n\theta)$ beyond standard textbook forms.
\item Examine what happens to the roots of a complex polynomial when coefficients are slightly perturbed from real to complex values.
\end{itemize}

\subsection*{Roots of Unity and Loci}
\begin{itemize}
\item Explore the locus of points satisfying $|z - 1| = |z + 1|$ and extend this idea to $|z - a| = |z - b|$ for complex constants $a$ and $b$.
\item Investigate how the sum of the $n$th roots of unity behaves when weighted by powers, such as $\sum_{k=0}^{n-1} k \omega^k$.
\item Study geometric transformations produced by multiplying a complex number by a fixed root of unity.
\end{itemize}

\subsection*{Hyperbolic Functions}
\begin{itemize}
\item Investigate the relationship between hyperbolic identities and circular trigonometric identities using exponential definitions.
\item Explore how $\cosh x$ and $\sinh x$ describe the geometry of a catenary curve and compare it with a parabola.
\item Study how inverse hyperbolic functions arise naturally in solving integrals involving quadratic expressions.
\end{itemize}

\subsection*{Differential Equations}
\begin{itemize}
\item Investigate how different initial conditions affect the family of solutions to $\frac{dy}{dx} = ky$ and interpret this as a growth model.
\item Explore the behaviour of solutions to $\frac{d^2y}{dx^2} + \omega^2 y = 0$ under varying $\omega$ and connect this to oscillatory motion.
\item Study how substitution methods can transform nonlinear differential equations into linear ones.
\end{itemize}

\subsection*{Polar Coordinates}
\begin{itemize}
\item Investigate how curves of the form $r = a + b \sin \theta$ change shape as $a$ and $b$ vary.
\item Explore conditions under which polar curves intersect themselves and how to detect these intersections analytically.
\item Study how converting between Cartesian and polar form changes the complexity of curve equations.
\end{itemize}

\subsection*{Matrices and Transformations}
\begin{itemize}
\item Investigate how repeated application of a matrix transformation affects geometric shapes in the plane.
\item Explore conditions under which a matrix represents a rotation, reflection, or shear.
\item Study eigenvalues as scaling factors and interpret them geometrically in two-dimensional transformations.
\end{itemize}

\subsection*{Further Investigation}
\begin{itemize}
\item Explore how different topics in this syllabus connect, such as complex numbers and rotations, or differential equations and exponential functions.
\item Investigate whether a single mathematical structure can unify transformations, differential equations, and complex numbers.
\end{itemize}

\vspace*{\fill}
\hrulefill
\begin{center}
\textit{“The beauty of mathematics only reveals itself to more patient followers.”}

Maryam Mirzakhani
\end{center}

\vspace*{\fill}
\end{document}